%% kapa: exam=1 round=33 year=2022 period=1 form=A subject=economics subject_ko=경제학원론 seq=33 number=73
%% source: Q-Net 감정평가사 기출문제 (cst00309) artlSeq=5209997 file=raw/1008/5209997/2022년도 제33회 감정평가사 1차시험 1교시(민법,경제학,부동산학).pdf
%% conversion: kordoc 4.24.51 -> tools/gen1.py (auto); review=proofread-sonnet; flags=none
\begin{question}{73}
경제성장모형인 Y=AK 모형에서 A는 0.5이고 저축률은 s, 감가상각률은 $\delta$일 때 이에 관한 설명으로 옳은 것은? (단, Y는 생산량, K는 자본량, $0<s<1$, $0<\delta<1$이다.)
\begin{choices}
  \item 자본의 한계생산은 체감한다.
  \item $\delta$ = 0.1이고 s = 0.4이면 경제는 지속적으로 성장한다.
  \item 감가상각률이 자본의 한계생산과 동일하면 경제는 지속적으로 성장한다.
  \item $\delta$ = s이면 경제는 균제상태(steady-state)이다.
  \item 자본의 한계생산이 자본의 평균생산보다 크다.
\end{choices}
\end{question}
\answer{2}
